\documentclass[10pt]{article}
\usepackage[margin=0.72in]{geometry}
\usepackage[T1]{fontenc}
\usepackage{lmodern}
\usepackage{amsmath,amssymb,amsthm,mathtools}
\usepackage{booktabs,longtable,array,enumitem,microtype,graphicx}
\usepackage[hidelinks]{hyperref}
\hypersetup{pdftitle={Technical Supplement: Exact Diffusion Design for Maximally Collective Stable Turing Patterns in Binary-Complex Mass-Action Networks},pdfauthor={Alec Kriebel}}
\newtheorem{theorem}{Theorem}[section]
\newtheorem{lemma}[theorem]{Lemma}
\newtheorem{proposition}[theorem]{Proposition}
\newcommand{\R}{\mathbb R}
\newcommand{\diag}{\operatorname{diag}}
\newcommand{\1}{\mathbf1}
\newcommand{\rank}{\operatorname{rank}}
\newcommand{\im}{\operatorname{im}}
\renewcommand{\thesection}{S\arabic{section}}
\makeatletter
\renewcommand*\l@section{\@dottedtocline{1}{1.5em}{2.8em}}
\makeatother
\allowdisplaybreaks
\title{Technical Supplement to\\[0.2em]
\large Exact Diffusion Design for Maximally Collective Stable Turing Patterns\\
\normalsize in Binary-Complex Mass-Action Networks}
\author{Alec Kriebel}
\date{August 2026}
\begin{document}
\maketitle
\tableofcontents

\section{Reaction matrices, flux cone, and semipositive conservation}

Order species as $X_1,\ldots,X_m,Z$ and reactions as
$R_0,R_2,\ldots,R_{m-2},R_a,R_b,R_+,R_-$.  For each indexed reaction, the
source and target vectors have total molecularity at most two.  The source
matrix $Y_m$ and stoichiometric matrix $\Gamma_m$ are obtained by placing
these source vectors and reaction differences in columns.

The balance equations $\Gamma_mv=0$ read, successively, equality of every
flux around the long circuit and equality of the two boundary-pair fluxes.
Thus
\[
 \ker\Gamma_m=
 \operatorname{span}\{(\1_m,0,0),(0_m,1,1)\}.
\]
An $m\times m$ minor formed by deleting one circuit reaction and one member of
the reversible pair is triangular after elementary row operations and has
absolute determinant $2^{m-2}$.  Therefore $\rank\Gamma_m=m$.
The vector
\[
 c=(0,4,\ldots,4,2,1)^T
\]
is orthogonal to every reaction vector, so $\im\Gamma_m=c^\perp$.
Notice that $c$ is semipositive rather than positive: $X_1$ is absent.

For flux $(a\1_m,b,b)$ and unit equilibrium,
$A_m(a,b)=\Gamma_m\diag(a\1_m,b,b)Y_m^T$ has precisely the entries printed in
the main text.  At an arbitrary positive equilibrium, right multiplication by
$H=\diag(1/x_i^*)$ gives $J=A_m(a,b)H$.  Conversely the rate assignment
$k_r=v_r/(x^*)^{y_r}$ realizes every $a,b,H>0$.

\section{All-spectrum principal-block classification}

Use the directed convention $X_j\to X_i$ when the $(i,j)$ Jacobian entry is
nonzero.  The interior arrows are the one-way chain
$X_2\to X_3\to\cdots\to X_{m-1}$.  The boundary arrows are those incident to
$X_1,X_m,Z$ and the returns $X_{m-1}\to X_1$, $X_m\to X_2$.  The following
case analysis is exhaustive for an induced set with fewer than $m$ vertices.

\begin{enumerate}[leftmargin=2em]
\item If a component contains a chain arrow and returns through
$X_{m-1}\to X_1$, then the unique forward path from $X_1$ forces every vertex
$X_1,\ldots,X_{m-1}$.  This is the first complete long cycle.
\item If it contains a chain arrow and returns through $X_m\to X_2$, then the
forward path forces every vertex $X_2,\ldots,X_m$.  This is the second complete
long cycle.
\item The alternative chain closure
$X_1\to X_2\to\cdots\to X_m\to X_1$ uses $m$ vertices and is excluded.
Adding any boundary vertex to either complete long cycle also uses at least
$m$ vertices.
\item If neither long cycle is complete, every feedback path not involving a
complete chain lies inside $\{X_1,X_m,Z\}$.  All other retained chain vertices
have no return path and therefore form negative singleton diagonal blocks.
\end{enumerate}
At $b=2a$, the coefficient $A_{m1}$ vanishes.  This deletes the arrow
$X_1\to X_m$; edge deletion can split but cannot merge strongly connected
components, so the same list remains exhaustive.

For the boundary triad, the Routh--Hurwitz difference expands as
\begin{align*}
c_1c_2-c_3={}&a\bigl(4 a^{2} h_{1}^{2} h_{m}+16 a^{2} h_{1} h_{m}^{2}+11 a b h_{1}^{2} h_{m}+4 a b h_{1}^{2} h_{Z}\\
&+32 a b h_{1} h_{m}^{2}+32 a b h_{1} h_{Z} h_{m}+64 a b h_{Z} h_{m}^{2}+7 b^{2} h_{1}^{2} h_{m}\\
&+4 b^{2} h_{1}^{2} h_{Z}+7 b^{2} h_{1} h_{m}^{2}+48 b^{2} h_{1} h_{Z} h_{m}+16 b^{2} h_{1} h_{Z}^{2}\\
&+16 b^{2} h_{Z} h_{m}^{2}+64 b^{2} h_{Z}^{2} h_{m}\bigr).
\end{align*}

Every monomial is strictly positive.  The two long-cycle modulus inequalities
are strict because $a+b>a$ and $4a+b>4a$.  Frobenius permutation therefore
proves Hurwitz stability of every smaller principal block.

Sparse expansion on $\{X_1,\ldots,X_m\}$ gives
\[
 (-1)^m\det J_{\widehat Z,\widehat Z}
 =-2a^{m-1}b\prod_{i=1}^m h_i.
\]
The complete omission table is derived in Section~S4 below.

\section{The principal-minor diffusion-ray theorem in full detail}


Let $a_I=(-1)^{|I|}\det J_{I,I}$ and $a_\varnothing=1$.  Assume
$\det J=0$, $a_I>0$ for $|I|\le n-2$, and
$\sum_{|I|=n-1}a_I>0$.  The last condition follows from a simple zero and
stable complementary spectrum but is stated directly because it is all the
proof uses.  Multilinearity in the diagonal entries of $sD$ gives
\[
 \det(sD-J)=\sum_{I\subseteq[n]}a_I
 \prod_{j\notin I}sd_j.
\]
Regrouping by $k=n-|I|$ yields
\[
 \det(sD-J)=\sum_{k=1}^n\beta_k(D)s^k,
 \qquad
 \beta_k(D)=\sum_{|I|=n-k}a_I\prod_{j\notin I}d_j.
\]
For $k\ge2$, every contributing set has order at most $n-2$, so the stated coefficient hypothesis gives $a_I>0$ and hence $\beta_k>0$.

The bracket $q_D(s)=\beta_1+\beta_2s+\cdots+\beta_ns^{n-1}$ is strictly
increasing for $s>0$.  If $\beta_1<0$, it has exactly one positive zero and
its derivative there is positive.  This proves scalar simplicity of the
dispersion root.  Ordinary algebraic simplicity of zero as an eigenvalue of
$J-s_*D$ follows from the characteristic-polynomial derivative below.

For the positive-real-eigenvalue band, differentiate
\[
 \chi_s(\lambda)=\sum_Ia_I\prod_{j\notin I}(\lambda+sd_j).
\]
Every order-at-most-$(n-2)$ contribution is nonnegative for $\lambda\ge0$.
The order-$(n-1)$ contribution to $\chi_s'$ is the constant
$\sum_{|I|=n-1}a_I$.  This is positive by hypothesis.  Under the common spectral corollary it equals the coefficient of $s$ in $\det(sI-J)$, and equal damping sends the simple zero and all stable eigenvalues left.  Thus $\chi_s'>0$ on the
nonnegative real axis.  The sign of $\chi_s(0)$ then determines whether there
is exactly one positive real eigenvalue.  At $s=s_*$,
$\chi_{s_*}(0)=0$ and $\chi_{s_*}'(0)>0$, so zero is an algebraically simple
eigenvalue of $J-s_*D$.

The theorem says nothing about nonreal unstable eigenvalues when $s>s_*$.
The all-mode certificates below are separate, profile-specific results.

\section{Order-\texorpdfstring{$(n-1)$}{(n-1)} omission minors and diffusion design}

Right multiplication by $H$ multiplies a principal determinant by the product
of its retained $h_i$, so the combinatorial coefficients may be computed at
$H=I$.

Omitting $Z$ leaves the complete $X$-chain block.  Its sparse cycle-cover
expansion gives
\[
 (-1)^m\det J_{\widehat Z,\widehat Z}
 =-2a^{m-1}b\prod_{i=1}^mh_i.
\]
For $2\le j\le m-1$, omit $X_j$ and order the remaining vertices by the two
feed-forward chain fragments, followed by $\{X_1,X_m,Z\}$.  The resulting
matrix is block triangular.  The $m-3$ chain singletons contribute signed
factors $a h_i$, and the boundary triad contributes
\[
 (-1)^3\det J_{\{X_1,X_m,Z\},\{X_1,X_m,Z\}}
 =16a^2b\,h_1h_mh_Z.
\]
Therefore
\[
 (-1)^m\det J_{\widehat{X_j},\widehat{X_j}}
 =16a^{m-1}bh_Z\prod_{\substack{1\le i\le m\\i\ne j}}h_i.
\]
Omitting $X_1$ leaves the restricted semipositive conservation vector in the
left kernel.  Omitting $X_m$ leaves the restriction of
$H^{-1}(2,-2,\ldots,-2,0,1)^T$ in the right kernel, because the omitted
coordinate is zero.  Hence both determinants vanish.

Multiplication by the omitted diffusivities produces
\[
 \beta_1(D)=2a^{m-1}b\left(\prod_{i=1}^mh_i\right)
 \left(8h_Z\sum_{j=2}^{m-1}\frac{d_j}{h_j}-d_Z\right).
\]
This yields the exact diffusion criterion and explains the coefficient eight
as the exact ratio $16/2$.

Define the stable realization domain
\[
 \mathfrak S_m=\{(a,b,H):a,b>0,\ H\succ0,\
 A_m(a,b)H|_{c^\perp}\text{ is Hurwitz}\}.
\]
The complete region $\mathfrak S_m$ is not characterized here.  For any
$(a,b,H)\in\mathfrak S_m$, if $d_{\min}=\min_i d_i$, then
\[
 d_Z>8h_Z\sum_{j=2}^{m-1}\frac{d_j}{h_j}
 \ge d_{\min}T(H),
\]
so $\chi_D>T(H)$.  Setting all $X$ diffusivities to one and
$d_Z=T(H)+\varepsilon$ proves the nonattained fixed-$H$ infimum.  Since
$h_Z/h_j\ge1/\chi_H$, one obtains $\chi_D\chi_H>8(m-2)$.  The implication
$T(H)>1$ is necessary under homogeneous stability; its converse is not proved.

\section{Improved unit-equilibrium critical profile}

Put $K_i=91m-181-i$.  The right vector and diffusion profile are
\[
 r_1=1,\quad r_i=-\frac{K_i}{63(m-2)},\quad
 r_m=-\frac29,\quad r_Z=\frac5{14},
\]
\[
 d_1=\frac{23}{63},\quad d_i=\frac1{K_i},\quad
 d_m=\frac17,\quad d_Z=\frac{16}{45}.
\]
Direct substitution gives $(A_m-D_m)r=0$.  With $\ell_Z=1$, the left vector is
\[
 \ell_1=-\frac{266}{815},\quad
 \ell_i=\frac{78260(m-2)}{163(91m-180-i)},\quad
 \ell_m=\frac{18368}{7335},\quad \ell_Z=1.
\]
Define
\[
 \mathfrak h_m=\sum_{j=1}^{m-2}\frac1{91m-181-j}.
\]
Then
\[
 \ell^Tr=-\frac{7043400m-13600927-7043400\mathfrak h_m}{924210}<0,
\]
\[
 \ell^TD_mr=-\frac{2(1760850\mathfrak h_m+16559)}{462105}<0.
\]
The harmonic bound $\mathfrak h_m\le(m-2)/(90m-179)$ makes the first sign an elementary
shifted-polynomial inequality.

\subsection*{Homogeneous and spatial determinants}
The homogeneous determinant has
\[
 \det(\lambda I-A_m)=\lambda q_m(\lambda),
 \qquad
 \lambda q_m=(1+\lambda)^{m-3}P(\lambda)-R(\lambda),
\]
with
$P=\lambda^4+12\lambda^3+42\lambda^2+47\lambda+16$ and
$R=5\lambda^2+33\lambda+16$.
For $m=3$,
\[
 q_3=(\lambda+7)(\lambda^2+5\lambda+2)
\]
is Hurwitz.  For $m\ge4$, the exact 35-term certificate for
$|1+\lambda|^2|P|^2-|R|^2$ is strictly positive away from the origin in the
closed right half-plane.  Since $|1+\lambda|\ge1$ there, a nonzero
closed-right-half-plane root of $q_m$ would contradict
$(1+\lambda)^{m-3}P=R$.  Finally, $q_m(0)=16m-34>0$, so the conservation zero
is simple and the complementary homogeneous spectrum is stable for every
$m\ge3$.
For a mode with damping factor $t$,
\[
 \det(\lambda I-A_m+tD_m)=Q_mF-G,
\]
where
\[
 Q_m=\prod_{i=2}^{m-1}\left(\lambda+1+\frac{t}{K_i}\right),
\]
\[
 g_1=\lambda+2+\frac{23}{63}t,\quad
 g_m=\lambda+5+\frac17t,\quad
 g_Z=\lambda+4+\frac{16}{45}t,
\]
\[
 F=g_1g_mg_Z-4g_1-4g_m+g_Z,\qquad
 G=g_Z(4g_1+g_m)-36.
\]
At onset, $\prod_{i=2}^{m-1}(1+K_i^{-1})=91/90$.  The exact certificate
tables appear in Section~S8.

\section{Conservation-compatible stable-mode corrections}

The quadratic tensor is obtained reaction by reaction.  Its components are
\begin{align*}
(B(u,v))_1&=-(u_1v_{m-1}+u_{m-1}v_1)+2u_Zv_Z-(u_1v_m+u_mv_1),\\
(B(u,v))_2&=-(u_1v_2+u_2v_1)+2u_mv_m,\\
(B(u,v))_i&=u_1v_{i-1}+u_{i-1}v_1-u_1v_i-u_iv_1,\\
(B(u,v))_m&=2(u_1v_{m-1}+u_{m-1}v_1)-4u_mv_m+2u_Zv_Z
 -(u_1v_m+u_mv_1),\\
(B(u,v))_Z&=-4u_Zv_Z+2(u_1v_m+u_mv_1).
\end{align*}
Here the third line holds for $3\le i\le m-1$.  Reaction-wise conservation
implies $c^TB(u,v)=0$.

The zero mode solves
\[
 A_mw_0=-\frac14B(r,r),\qquad c^Tw_0=0.
\]
The exact solution is
\[
(w_0)_1=\frac{182448m-373417}{31752(8m-17)},
\]
\[
(w_0)_2=\frac{1008m^2-20459m+37138}{31752(m-2)(8m-17)},
\]
\[
(w_0)_i=(w_0)_2-\frac{i-2}{126(m-2)},\quad
(w_0)_m=-\frac1{81},\quad
(w_0)_Z=\frac{16861m-34044}{7938(8m-17)}.
\]

The second harmonic solves
\[
 (A_m-4D_m)w_2=-\frac14B(r,r).
\]
With
\[
 \sigma=\frac1{126(m-2)},\qquad
 T_i=\frac{K_{i-3}K_{i-2}K_{i-1}K_i}{K_{-1}K_0K_1K_2},
\]
the chain recurrence is
\[
 (w_2)_i=T_i\left((w_2)_2+\frac{\sigma K_2}{3}\right)
 -\frac{\sigma K_i}{3}.
\]
Its boundary denominator is
\[
 \mathcal Q_m=589180301m^3-3500015940m^2+6930529579m-4574434500.
\]
After $m=u+3$, the coefficients are
\[
589180301,\quad1802606769,\quad1838302066,\quad624878904,
\]
so it never vanishes for $m\ge3$.  The complete boundary values are provided
in the machine-readable certificate and in the public verifier.

The cubic coefficient is
\[
 c_m=\frac{\ell^T\{B(r,w_0)+\tfrac12B(r,w_2)\}}{\ell^Tr}.
\]
This is the cubic coefficient of the one-dimensional center-manifold normal
form in the adjoint-projection coordinate.  Setting the reduced vector field
to zero gives the associated stationary Lyapunov--Schmidt equation.
For the cubic numerator, introduce
\begin{align*}
 P_R(m)={}&68605040480814208768m^4
 -550882186169626030957m^3\\
 &+1658612632937449670852m^2
 -2219226476204103501323m\\
 &+1113379274975809565700,
\end{align*}
and
\begin{align*}
 P_C(m)={}&652054120726848m^4
 -5151971981328467m^3\\
 &+15265080924982572m^2
 -20102347725659113m\\
 &+9927281930180400.
\end{align*}
Then the exact decomposition is
\[
 N_m=N_m^{\rm ref}=R_m+C_m\mathfrak h_m,
\]
where
\[
 R_m=\frac{P_R(m)}
 {286118780220(8m-17)\mathcal Q_m},
 \qquad
 C_m=-\frac{215P_C(m)}
 {11645046(8m-17)\mathcal Q_m}.
\]
The shifted-positive coefficient certificates in Section~S8.2 show that
$\mathcal Q_m>0$, $P_R(m)>0$, and $P_C(m)>0$ for $m\ge3$.
Consequently $R_m>0$ and $C_m<0$; in the latter conclusion the external
factor $-215$ reverses the sign of the positive polynomial $P_C$.

The exact clearing identity is
\[
 R_m+C_m\frac{m-2}{90m-179}
 =\frac{L_m}
 {286118780220(8m-17)(90m-179)\mathcal Q_m},
\]
where
\begin{align*}
L_m={}&2729945147827667886720m^5
-27755132420474170999952m^4\\
&+112813395868533457497683m^3
-229153280695458887386228m^2\\
&+232620996871721820873517m
-94412163900120968220300.
\tag{S6.1}\label{eq:Lmpoly}
\end{align*}
After $m=u+3$, its coefficients are
\begin{gather*}
2729945147827667886720,
13194044796940847300848,
25446870127333515303059,\\
24475321329207325509911,
11736484608366875120374,
2243933770033305838488,
\end{gather*}
all positive.  Since $C_m<0$ and
$\mathfrak h_m\le(m-2)/(90m-179)$, the complete comparison is
\[
 R_m+C_m\mathfrak h_m
 \ge R_m+C_m\frac{m-2}{90m-179}>0.
\]
Thus the cubic numerator is positive and, since $\ell^Tr<0$, $c_m<0$.

\section{Equilibrium-scaled stable family}

Let $\nu=m-2$ and define
\[
 \kappa_\nu=\begin{cases}
 1/\sqrt3,&\nu=1,\\
 \sqrt5/2,&\nu\ge2,
 \end{cases}
 \qquad
 L_0(\nu)=\frac{\kappa_\nu}{\sqrt\nu},
 \qquad
 L_1(\nu)=\frac{90\nu}{90\nu+1}.
\]
This interval is nonempty.  For $\nu=1$, square both positive sides and use
$8281<24300$.  For $\nu\ge2$,
\[
 (90\nu+1)^2\le(91\nu)^2<6480\nu^3,
\]
which is equivalent to $L_0(\nu)<L_1(\nu)$.  The lower endpoints are
sufficient boundaries of the certificates below, not proved intrinsic
dynamical boundaries.  For $L\in[L_0,L_1]$, set
\[
 h_1=h_m=h_Z=1,\qquad h_i=\frac{K_i}{LK_{i-1}}.
\]
Write $\mathsf H_m(L)=\diag(h_1,\ldots,h_m,h_Z)$.  The physical
equilibrium is $x^*=\mathsf H_m(L)^{-1}\1$ and the physical diffusion matrix
is $D^{\rm phys}=\mathsf H_m(L)\Delta$, where $\Delta$ is the improved unit
profile.  In normalized variables $z=\mathsf H_m(L)x$,
\[
 z_t=\mathsf H_m(L)\{f_m(z)+\Delta z_{\xi\xi}\}.
\]
The stationary bracket and critical right vector are unchanged.  The dynamic
left vector is
\[
 \widetilde\ell_m(L)=\mathsf H_m(L)^{-1}\ell_m,
\]
and the physical conservation gauge uses $\mathsf H_m(L)^{-1}c$.

For interior chain factors,
\[
 g_i=\frac{K_{i-1}}{K_i}(1+L\lambda)+\frac{t-1}{K_i}.
\]
When $t\ge1$ and $\Re\lambda\ge0$,
\[
 |Q_m|^2\ge\left(\frac{91}{90}\right)^2|1+L\lambda|^{2\nu}
 \ge\left(\frac{91}{90}\right)^2
 \left(1+2\nu L\Re\lambda+\frac13(\Im\lambda)^2\right).
\]
The last inequality holds because $\nu L^2\ge1/3$ throughout the certified
interval.  The 84-term spatial coefficient table in Section~S8 is
nonnegative and is strict except at $(\lambda,t)=(0,1)$.

Homogeneous stability uses a different boundary polynomial.  At $t=0$,
\[
 \frac{\det(\lambda I-\mathsf H_m(L) A_m)}{\det \mathsf H_m(L)}
 =Q(\lambda)F_0(\lambda)-R(\lambda),
\]
where
\[
 Q=\prod_{i=2}^{m-1}\left(1+L\frac{K_{i-1}}{K_i}\lambda\right),\qquad
 F_0=\lambda^3+11\lambda^2+31\lambda+16,
 \qquad R=5\lambda^2+33\lambda+16.
\]
For $\Re\lambda=x\ge0$, every factor in $Q$ has modulus at least
$|1+L\lambda|$.  If $\nu\ge2$, then
\[
 |Q|^2\ge |1+L\lambda|^{2\nu}
 \ge 1+2\nu Lx+\nu L^2(\Im\lambda)^2
 \ge 1+Ax+\frac54(\Im\lambda)^2,
 \qquad A=2\nu L.
\]
Here $A\ge\sqrt{5\nu}\ge\sqrt{10}>1/4$.
The exact polynomial
\[
 \left(1+Ax+\frac54(\Im\lambda)^2\right)|F_0|^2-|R|^2
\]
has 22 nonzero monomials.  With $U=A-1/4$, its coefficient polynomials have
nonnegative rational coefficients and the expression is strict away from
$\lambda=0$.  Thus there is no nonzero homogeneous eigenvalue in the closed
right half-plane.  The coefficient $5/4$ is sharp for this coefficientwise
certificate.  If it
is replaced by $B$, then at $x=0$ the coefficient of $z=y^2$ in
$(1+Bz)|F_0|^2-|R|^2$ is $256B-320$.  It vanishes at $B=5/4$ and is negative
for $B<5/4$.  This does not show that $L_0$ is an intrinsic dynamical
boundary; it identifies only the boundary of the present exact modulus
certificate.  Moreover
\[
 (QF_0-R)'(0)=16L\sum_{i=2}^{m-1}\frac{K_{i-1}}{K_i}-2>0,
\]
where the inequality uses $K_{i-1}/K_i>1$ and
$L\ge\sqrt5/(2\sqrt\nu)$.  Thus the conservation zero is simple.  If
$\nu=1$, put $c=91L/90$.  Direct
division by the conservation factor gives
\[
 \frac{QF_0-R}{\lambda}
 =c\lambda^3+(1+11c)\lambda^2+(6+31c)\lambda+(16c-2).
\]
Here $c>1/8$, and the remaining cubic Routh--Hurwitz difference is
$6+99c+325c^2>0$.  This proves homogeneous relative stability in the
exceptional dimension as well.

Let
$\rho_m=(2,-2,\ldots,-2,0,1)^T$.  The physical zero-mode gauge is
\[
 w_0(L)=w_0^{\rm ref}+\tau_m(L)\rho_m,
 \qquad
 \tau_m(L)=-\frac{(\mathsf H_m(L)^{-1}c)^Tw_0^{\rm ref}}
 {(\mathsf H_m(L)^{-1}c)^T\rho_m}.
\]
The second harmonic is unchanged.  The cubic numerator satisfies
\[
 N_m(L)=N_m^{\rm ref}+\tau_m(L)S_m,
\]
where $N_m^{\rm ref}=R_m+C_m\mathfrak h_m$.  To make its quantitative margin
explicit, set
\begin{align*}
 P_{\rm ref}(\nu)={}&3790502986637265684840\nu^5
 -974216530468600286489\nu^4\\
 &-53103567440921218871\nu^3
 -576386186827093561\nu^2\\
 &+3649732858601219\nu+55281268032918
\end{align*}
and
\begin{align*}
 D_{\rm ref}(\nu)={}&715296950550(8\nu-1)(90\nu+1)\\
 &\mathrel{}\times
 (589180301\nu^3+35065866\nu^2+629431\nu+3306).
\end{align*}
Since $C_m<0$ and $\mathfrak h_m\le\nu/(90\nu+1)$, exact clearing gives
\[
 N_m^{\rm ref}-\frac1{100}
 \ge \frac{P_{\rm ref}(\nu)}{D_{\rm ref}(\nu)}.
\]
The denominator is positive for $\nu\ge1$.  After $\nu=v+1$, the coefficients
of $P_{\rm ref}$, in descending powers $v^5,\ldots,v^0$, are
\begin{gather*}
3790502986637265684840,
17978298402717728137711,
33955060177057334483573,\\
31899843595051464379292,
14895188986348368035728,
2762610207555043720056.
\end{gather*}
They are positive, so $P_{\rm ref}(\nu)>0$ and
$N_m^{\rm ref}>1/100$.

Separately,
\[
 S_m=-\frac{4(1760850\mathfrak h_m-10253)}{462105}.
\]
The bounds $1/91\le\mathfrak h_m<1/90$ give $S_m<0$ because
\[
 \frac{1760850}{91}-10253>0,
\]
and give $S_m>-1/10$ because
\[
 \frac4{462105}\left(\frac{1760850}{90}-10253\right)<\frac1{10}.
\]

For the remaining gauge estimate, define
\begin{align*}
 A_\tau={}&1494249120\mathfrak h_mL\nu^2
 -69786990\mathfrak h_mL\nu+108738630L\nu^2\\
 &+1214388L\nu-8521L-125249670\nu^2+1031940\nu,
\end{align*}
\[
 B_\tau=32760\mathfrak h_mL\nu+32760L\nu^2+4L-4095\nu.
\]
Then the closed form is
\[
 \tau_m(L)=-\frac{A_\tau}{15876(8\nu-1)B_\tau}.
\]
On the broader certified cubic range $L\ge1/\sqrt{3\nu}$,
\[
 B_\tau=4095\nu\bigl[8L(\nu+\mathfrak h_m)-1\bigr]+4L>0,
\]
because $L\nu\ge1/\sqrt3$.  Exact differentiation then gives
\[
 \partial_{\mathfrak h}\tau_m
 =-\frac{4225L\nu^2(182448L\nu+1008L-7513)}{2B_\tau^2}<0,
\]
\[
 \partial_L\tau_m
 =-\frac{65\nu(-61531470\mathfrak h_m\nu
 +125249670\nu^2+1031940\nu-7513)}{252B_\tau^2}<0.
\]
For the first sign, $L\nu\ge1/\sqrt3>1/2$ makes the parenthesis positive.
For the second, use $\mathfrak h_m<1/90$ and $\nu\ge1$.  Thus the maximum
occurs at $\mathfrak h_m=1/91$ and $L=1/\sqrt{3\nu}$.

Put $y=\sqrt{3\nu}$ and define the unreduced endpoint polynomial
\begin{align*}
 P_\tau(\nu,y)={}&-1040195520\nu^3+756272790\nu^2y
 -507201030\nu^2\\
 &-21412755\nu y-935658\nu+58481
\end{align*}
and $D_\tau(\nu,y)=-32760\nu^2+4095\nu y-360\nu-4$.  Direct
substitution gives
\[
 \tau_m(1/91,1/y)-\frac1{20}
 =-\frac{P_\tau(\nu,y)}{79380(8\nu-1)D_\tau(\nu,y)}.
\]
Since $y<2\nu$, $D_\tau<0$.  For $\nu=1$, the coefficient of $y$ in
$P_\tau(1,y)$ is positive and $\sqrt3<7/4$, whence
\[
 P_\tau(1,\sqrt3)<P_\tau(1,7/4)=-\frac{1049074663}{4}<0.
\]
For $\nu\ge2$, $y\le5\nu/4$ and deletion of the negative term
$-21412755\nu y$ gives $P_\tau(\nu,y)<P_{\rm up}(\nu)$, where
\[
 P_{\rm up}(\nu)=-\frac{189709065}{2}\nu^3-507201030\nu^2
 -935658\nu+58481.
\]
This upper-bound polynomial is human-checkable:
\[
 P_{\rm up}'(\nu)=-\frac{569127195}{2}\nu^2-1014402060\nu-935658<0
 \qquad(\nu\ge0),
\]
and
\[
 P_{\rm up}(2)=-2789453215<0.
\]
Hence $P_{\rm up}(\nu)<0$ for every $\nu\ge2$, and therefore
$P_\tau(\nu,y)<0$ in all dimensions.  Because the displayed denominator
$79380(8\nu-1)D_\tau$ is negative, the exact endpoint comparison proves
$\tau_m(L)<1/20$.

If $\tau_m(L)\le0$, then $\tau_m(L)S_m\ge0$.  If instead
$0<\tau_m(L)<1/20$, then $\tau_m(L)S_m>-1/200$.  Together with
$N_m^{\rm ref}>1/100$, this yields $N_m(L)>1/200$.  Also
\[
 \ell^T\mathsf H_m(L)^{-1}r=-\frac{485873}{924210}
 -\frac{11180}{1467}L(m-2)<0.
\]
Therefore the family is uniformly supercritical over the certified interval.
Indeed, the cubic comparison was proved on the broader range
$L\ge1/\sqrt{3\nu}$; the present lower endpoint agrees with this value for
$\nu=1$ and is larger for $\nu\ge2$.

The contrasts are
\[
 \chi_D(L)=\frac{23}{63}(91\nu L),\qquad
 \chi_H(L)=\frac{91\nu-1}{91\nu L},
\]
where $\chi_H$ measures equilibrium concentration-scale separation, not a
universal measure of biological expense or implausibility.  Their product is
fixed:
\[
 \chi_D(L)\chi_H(L)=\frac{23(91\nu-1)}{63}
 =\frac{23(91m-183)}{63}=\chi_D^{\rm unit}(m).
\]
Thus $L$ reallocates contrast between diffusion and equilibrium concentration
scales and reduces $\max(\chi_D,\chi_H)$, not their product.  At $L=L_0$,
\[
 \chi_D=\frac{2093}{63}\kappa_\nu\sqrt\nu,
 \qquad
 \chi_H=\frac{\sqrt\nu}{\kappa_\nu}\frac{91\nu-1}{91\nu},
\]
so both contrasts have square-root growth.

\section{Exact coefficient tables}\label{sec:coefftables}

\subsection{Coefficientwise-nonnegative modulus certificates}

The first four tables below are coefficientwise-nonnegative modulus
certificates.  Each row records an exponent vector and the exact coefficient
of the corresponding polynomial.  The later shifted-polynomial tables are
signed scalar certificates; for each such table we state the polynomial, any
overall sign prefactor, the parameter domain, and the sign conclusion
separately.

By conjugation, take $y\ge0$ and use the variables
\[
 \lambda=x+i y=x+i\sqrt z,\qquad z=y^2,\qquad t=1+s,
 \qquad x,y,z,s\ge0.
\]
For the unit-profile spatial calculation, define
\[
 g_1=\lambda+2+\frac{23}{63}t,\qquad
 g_m=\lambda+5+\frac17t,\qquad
 g_Z=\lambda+4+\frac{16}{45}t,
\]
\[
 F=g_1g_mg_Z-4g_1-4g_m+g_Z,
 \qquad G=g_Z(4g_1+g_m)-36.
\]
The 77-term table is exactly the expansion in $(x,z,s)$ of
\[
 E_{77}(x,z,s)=
 \left[\left(\frac{91}{90}\right)^2+z\right]
 \left|F(x+i\sqrt z,1+s)\right|^2
 -\left|G(x+i\sqrt z,1+s)\right|^2.
\]
Every one of its 77 displayed coefficients is strictly positive; there is no
constant term.  Hence equality occurs only at $x=z=s=0$, equivalently
$(\lambda,t)=(0,1)$.

For the equilibrium-scaled spatial calculation, put $A=2\nu L>0$.  The
84-term table is exactly the expansion in $(x,z,s)$ of
\[
 E_{84}(x,z,s;A)=
 \left(\frac{91}{90}\right)^2
 \left(1+Ax+\frac z3\right)
 \left|F(x+i\sqrt z,1+s)\right|^2
 -\left|G(x+i\sqrt z,1+s)\right|^2.
\]
Its 84 coefficients are polynomials in $A$ with nonnegative rational
coefficients and are positive for $A>0$.  Again there is no constant term,
and equality occurs only at $x=z=s=0$, equivalently
$(\lambda,t)=(0,1)$.

For the equilibrium-scaled homogeneous calculation, define
\[
 F_0(\lambda)=\lambda^3+11\lambda^2+31\lambda+16,
 \qquad R(\lambda)=5\lambda^2+33\lambda+16,
 \qquad U=A-\frac14.
\]
The 22-term table is exactly the expansion in $(x,z)$ of
\[
 E_{22}(x,z;U)=
 \left[1+\left(U+\frac14\right)x+\frac54z\right]
 |F_0(x+i\sqrt z)|^2-|R(x+i\sqrt z)|^2.
\]
For $U\ge0$, all 22 coefficient polynomials are nonnegative, and the
expression vanishes only at $x=z=0$, equivalently $\lambda=0$.

Finally, with
\[
 P(\lambda)=\lambda^4+12\lambda^3+42\lambda^2+47\lambda+16,
\]
the unit-profile homogeneous table is the 35-term expansion in $(x,z)$ of
\[
 E_{35}(x,z)=|1+\lambda|^2|P(\lambda)|^2-|R(\lambda)|^2,
 \qquad \lambda=x+i\sqrt z.
\]
All 35 coefficients are strictly positive, there is no constant term, and
equality occurs only at $x=z=0$, equivalently $\lambda=0$.

From the repository root, the command
\texttt{python independent\_verifier/verify\_symbolic\_certificates.py}
regenerates and checks the tables from exact symbolic expressions.

\subsubsection*{35-term homogeneous certificate}
\begin{longtable}{rrl}
$\deg_{x}$ & $\deg_{z}$ & coefficient\\
\toprule
\endfirsthead
$\deg_{x}$ & $\deg_{z}$ & coefficient\\
\toprule
\endhead
10 & 0 & $1$\\
9 & 0 & $26$\\
8 & 1 & $5$\\
8 & 0 & $277$\\
7 & 1 & $104$\\
7 & 0 & $1582$\\
6 & 2 & $10$\\
6 & 1 & $892$\\
6 & 0 & $5356$\\
5 & 2 & $156$\\
5 & 1 & $4034$\\
5 & 0 & $11282$\\
4 & 3 & $10$\\
4 & 2 & $1014$\\
4 & 1 & $10432$\\
4 & 0 & $15116$\\
3 & 3 & $104$\\
3 & 2 & $3322$\\
3 & 1 & $15628$\\
3 & 0 & $12612$\\
2 & 4 & $5$\\
2 & 3 & $460$\\
2 & 2 & $5804$\\
2 & 1 & $13296$\\
2 & 0 & $5568$\\
1 & 4 & $26$\\
1 & 3 & $870$\\
1 & 2 & $4858$\\
1 & 1 & $5700$\\
1 & 0 & $960$\\
0 & 5 & $1$\\
0 & 4 & $61$\\
0 & 3 & $728$\\
0 & 2 & $1508$\\
0 & 1 & $192$\\
\bottomrule
\end{longtable}

\subsubsection*{77-term improved-profile spatial certificate}
\begin{longtable}{rrrl}
$\deg_{x}$ & $\deg_{z}$ & $\deg_{s}$ & coefficient\\
\toprule
\endfirsthead
$\deg_{x}$ & $\deg_{z}$ & $\deg_{s}$ & coefficient\\
\toprule
\endhead
6 & 1 & 0 & $1$\\
6 & 0 & 0 & $8281/8100$\\
5 & 1 & 1 & $544/315$\\
5 & 1 & 0 & $7474/315$\\
5 & 0 & 1 & $160888/91125$\\
5 & 0 & 0 & $4420871/182250$\\
4 & 2 & 0 & $3$\\
4 & 1 & 2 & $40058/33075$\\
4 & 1 & 1 & $1147126/33075$\\
4 & 1 & 0 & $9681647/44100$\\
4 & 0 & 2 & $3384901/2733750$\\
4 & 0 & 1 & $96932147/2733750$\\
4 & 0 & 0 & $1073317517/5467500$\\
3 & 2 & 1 & $1088/315$\\
3 & 2 & 0 & $14948/315$\\
3 & 1 & 3 & $2744576/6251175$\\
3 & 1 & 2 & $41276762/2083725$\\
3 & 1 & 1 & $8186375738/31255875$\\
3 & 1 & 0 & $10434543841/10418625$\\
3 & 0 & 3 & $115958336/258339375$\\
3 & 0 & 2 & $3487886389/172226250$\\
3 & 0 & 1 & $39671949511/172226250$\\
3 & 0 & 0 & $315432259813/516678750$\\
2 & 3 & 0 & $3$\\
2 & 2 & 2 & $147968/99225$\\
2 & 2 & 1 & $4065856/99225$\\
2 & 2 & 0 & $112938659/396900$\\
2 & 1 & 4 & $33947257/393824025$\\
2 & 1 & 3 & $2164641406/393824025$\\
2 & 1 & 2 & $1130666508023/9845600625$\\
2 & 1 & 1 & $9041228678666/9845600625$\\
2 & 1 & 0 & $5074377793279/2187911250$\\
2 & 0 & 4 & $5737086433/65101522500$\\
2 & 0 & 3 & $182912198807/32550761250$\\
2 & 0 & 2 & $2142193804481/21700507500$\\
2 & 0 & 1 & $17194597630733/32550761250$\\
2 & 0 & 0 & $34399269232129/65101522500$\\
1 & 3 & 1 & $544/315$\\
1 & 3 & 0 & $7474/315$\\
1 & 2 & 3 & $1817216/6251175$\\
1 & 2 & 2 & $24359162/2083725$\\
1 & 2 & 1 & $5056135954/31255875$\\
1 & 2 & 0 & $47437544293/62511750$\\
1 & 1 & 5 & $3399584/393824025$\\
1 & 1 & 4 & $57943048/78764805$\\
1 & 1 & 3 & $1934522477168/88610405625$\\
1 & 1 & 2 & $16308922895827/59073603750$\\
1 & 1 & 1 & $85574440190593/59073603750$\\
1 & 1 & 0 & $62146993882837/25317258750$\\
1 & 0 & 5 & $143632424/16275380625$\\
1 & 0 & 4 & $2448093778/3255076125$\\
1 & 0 & 3 & $98527787884/5425126875$\\
1 & 0 & 2 & $475980699652/3255076125$\\
1 & 0 & 1 & $4488365319374/16275380625$\\
1 & 0 & 0 & $1086412028/10333575$\\
0 & 4 & 0 & $1$\\
0 & 3 & 2 & $27794/99225$\\
0 & 3 & 1 & $624478/99225$\\
0 & 3 & 0 & $1747307/26460$\\
0 & 2 & 4 & $1744613/78764805$\\
0 & 2 & 3 & $437610814/393824025$\\
0 & 2 & 2 & $452220491743/19691201250$\\
0 & 2 & 1 & $4035916860091/19691201250$\\
0 & 2 & 0 & $10861695961357/13127467500$\\
0 & 1 & 6 & $135424/393824025$\\
0 & 1 & 5 & $2964608/78764805$\\
0 & 1 & 4 & $962538972893/637994920500$\\
0 & 1 & 3 & $43240103983367/1594987301250$\\
0 & 1 & 2 & $238300444324673/1063324867500$\\
0 & 1 & 1 & $170939874507539/227855328750$\\
0 & 1 & 0 & $32638690751329/65101522500$\\
0 & 0 & 6 & $5721664/16275380625$\\
0 & 0 & 5 & $125254688/3255076125$\\
0 & 0 & 4 & $6479423564/5425126875$\\
0 & 0 & 3 & $41256502232/3255076125$\\
0 & 0 & 2 & $513831327844/16275380625$\\
0 & 0 & 1 & $148103696/10333575$\\
\bottomrule
\end{longtable}

\subsubsection*{22-term equilibrium-scaled homogeneous certificate ($U=A-1/4$)}
\begin{longtable}{rrl}
$\deg_{x}$ & $\deg_{z}$ & coefficient\\
\toprule
\endfirsthead
$\deg_{x}$ & $\deg_{z}$ & coefficient\\
\toprule
\endhead
7 & 0 & $1/4+U$\\
6 & 1 & $5/4$\\
6 & 0 & $13/2+22U$\\
5 & 1 & $113/4+3U$\\
5 & 0 & $271/4+183U$\\
4 & 2 & $15/4$\\
4 & 1 & $971/4+44U$\\
4 & 0 & $673/2+714U$\\
3 & 2 & $223/4+3U$\\
3 & 1 & $997+242U$\\
3 & 0 & $2849/4+1313U$\\
2 & 3 & $15/4$\\
2 & 2 & $311+22U$\\
2 & 1 & $7919/4+586U$\\
2 & 0 & $312+992U$\\
1 & 3 & $111/4+U$\\
1 & 2 & $3077/4+59U$\\
1 & 1 & $6593/4+609U$\\
1 & 0 & $256U$\\
0 & 4 & $5/4$\\
0 & 3 & $299/4$\\
0 & 2 & $3181/4$\\
\bottomrule
\end{longtable}

\subsubsection*{84-term equilibrium-scaled spatial certificate}
\begin{longtable}{rrrl}
$\deg_{x}$ & $\deg_{z}$ & $\deg_{s}$ & coefficient\\
\toprule
\endfirsthead
$\deg_{x}$ & $\deg_{z}$ & $\deg_{s}$ & coefficient\\
\toprule
\endhead
7 & 0 & 0 & $8281A/8100$\\
6 & 1 & 0 & $8281/24300$\\
6 & 0 & 1 & $160888A/91125$\\
6 & 0 & 0 & $8281/8100+4420871A/182250$\\
5 & 1 & 1 & $160888/273375$\\
5 & 1 & 0 & $4420871/546750+8281A/2700$\\
5 & 0 & 2 & $3384901A/2733750$\\
5 & 0 & 1 & $160888/91125+96932147A/2733750$\\
5 & 0 & 0 & $4420871/182250+1210005017A/5467500$\\
4 & 2 & 0 & $8281/8100$\\
4 & 1 & 2 & $3384901/8201250$\\
4 & 1 & 1 & $96932147/8201250+321776A/91125$\\
4 & 1 & 0 & $315078023/4100625+4420871A/91125$\\
4 & 0 & 3 & $115958336A/258339375$\\
4 & 0 & 2 & $3384901/2733750+3487886389A/172226250$\\
4 & 0 & 1 & $96932147/2733750+45494837011A/172226250$\\
4 & 0 & 0 & $1073317517/5467500+503404909813A/516678750$\\
3 & 2 & 1 & $321776/273375$\\
3 & 2 & 0 & $4420871/273375+8281A/2700$\\
3 & 1 & 3 & $115958336/775018125$\\
3 & 1 & 2 & $3487886389/516678750+6251648A/4100625$\\
3 & 1 & 1 & $47319306931/516678750+171782416A/4100625$\\
3 & 1 & 0 & $578603925523/1550036250+2360113561A/8201250$\\
3 & 0 & 4 & $5737086433A/65101522500$\\
3 & 0 & 3 & $115958336/258339375+182912198807A/32550761250$\\
3 & 0 & 2 & $3487886389/172226250+2513945531981A/21700507500$\\
3 & 0 & 1 & $39671949511/172226250+29165273350733A/32550761250$\\
3 & 0 & 0 & $315432259813/516678750+138537167092129A/65101522500$\\
2 & 3 & 0 & $8281/8100$\\
2 & 2 & 2 & $6251648/12301875$\\
2 & 2 & 1 & $171782416/12301875+160888A/91125$\\
2 & 2 & 0 & $4871148347/49207500+4420871A/182250$\\
2 & 1 & 4 & $5737086433/195304567500$\\
2 & 1 & 3 & $182912198807/97652283750+76777376A/258339375$\\
2 & 1 & 2 & $2613196695629/65101522500+2058349189A/172226250$\\
2 & 1 & 1 & $33256099805357/97652283750+28172026051A/172226250$\\
2 & 1 & 0 & $184975683058783/195304567500+388034017813A/516678750$\\
2 & 0 & 5 & $143632424A/16275380625$\\
2 & 0 & 4 & $5737086433/65101522500+2448093778A/3255076125$\\
2 & 0 & 3 & $182912198807/32550761250+119438423884A/5425126875$\\
2 & 0 & 2 & $2142193804481/21700507500+878963880652A/3255076125$\\
2 & 0 & 1 & $17194597630733/32550761250+21944279796374A/16275380625$\\
2 & 0 & 0 & $34399269232129/65101522500+21842305628A/10333575$\\
1 & 3 & 1 & $160888/273375$\\
1 & 3 & 0 & $4420871/546750+8281A/8100$\\
1 & 2 & 3 & $76777376/775018125$\\
1 & 2 & 2 & $2058349189/516678750+2348593A/8201250$\\
1 & 2 & 1 & $29084261011/516678750+52768391A/8201250$\\
1 & 2 & 0 & $212816762834/775018125+1090212071A/16402500$\\
1 & 1 & 5 & $143632424/48826141875$\\
1 & 1 & 4 & $2448093778/9765228375+294839597A/13020304500$\\
1 & 1 & 3 & $124275398572/16275380625+36978113783A/32550761250$\\
1 & 1 & 2 & $9956722796683/97652283750+503148512477A/21700507500$\\
1 & 1 & 1 & $11312104230233/19530456750+6606581721077A/32550761250$\\
1 & 1 & 0 & $241757054977/221433750+52308962411329A/65101522500$\\
1 & 0 & 6 & $5721664A/16275380625$\\
1 & 0 & 5 & $143632424/16275380625+125254688A/3255076125$\\
1 & 0 & 4 & $2448093778/3255076125+8242156364A/5425126875$\\
1 & 0 & 3 & $98527787884/5425126875+86436565592A/3255076125$\\
1 & 0 & 2 & $475980699652/3255076125+3438519744244A/16275380625$\\
1 & 0 & 1 & $4488365319374/16275380625+7094103536A/10333575$\\
1 & 0 & 0 & $1086412028/10333575+4999696A/6561$\\
0 & 4 & 0 & $8281/24300$\\
0 & 3 & 2 & $2348593/24603750$\\
0 & 3 & 1 & $52768391/24603750$\\
0 & 3 & 0 & $570259573/24603750$\\
0 & 2 & 4 & $294839597/39060913500$\\
0 & 2 & 3 & $36978113783/97652283750$\\
0 & 2 & 2 & $521791643711/65101522500$\\
0 & 2 & 1 & $3617447476357/48826141875$\\
0 & 2 & 0 & $30203751676613/97652283750$\\
0 & 1 & 6 & $5721664/48826141875$\\
0 & 1 & 5 & $125254688/9765228375$\\
0 & 1 & 4 & $11480941147/21700507500$\\
0 & 1 & 3 & $975299997269/97652283750$\\
0 & 1 & 2 & $17163191001769/195304567500$\\
0 & 1 & 1 & $10010796373877/32550761250$\\
0 & 1 & 0 & $650029271329/65101522500$\\
0 & 0 & 6 & $5721664/16275380625$\\
0 & 0 & 5 & $125254688/3255076125$\\
0 & 0 & 4 & $6479423564/5425126875$\\
0 & 0 & 3 & $41256502232/3255076125$\\
0 & 0 & 2 & $513831327844/16275380625$\\
0 & 0 & 1 & $148103696/10333575$\\
\bottomrule
\end{longtable}


\subsection{Signed scalar and rational-function certificates}

Each generated row below states the exact scalar polynomial or rational
expression, the descending coefficient order after its indicated shift, the
sign of the polynomial on its parameter domain, and whether an external
negative prefactor reverses that sign.  The shifts use $u=m-3$ and
$v=\nu-1$, leaving $z=y^2$ reserved for the modulus tables.  In particular,
the $P_C$ row certifies $C_m<0$ through its external factor $-215$; the sign
of $S_m$ is certified separately.

\paragraph{Unit-profile boundary denominator.} After $m=z+3$, the coefficients of the cubic denominator are
\[
\begin{gathered}
589180301,\quad 1802606769,\quad 1838302066,\\
624878904\\
\end{gathered}
\]

\paragraph{Unit-profile critical-denominator bound.} After $m=z+3$, the auxiliary quadratic has coefficients
\[
\begin{gathered}
633906000,\quad 1311540570,\quad 678120443\\
\end{gathered}
\]

\paragraph{Unit-profile cubic numerator lower bound.} After $m=z+3$, the decisive degree-five polynomial has coefficients
\[
\begin{gathered}
2729945147827667886720,\quad 13194044796940847300848,\quad 25446870127333515303059,\\
24475321329207325509911,\quad 11736484608366875120374,\quad 2243933770033305838488\\
\end{gathered}
\]

\paragraph{Equilibrium-scaled coefficient sign.} After $m=z+3$, the polynomial whose sign gives $S_m<0$ has coefficients
\[
\begin{gathered}
652054120726848,\quad 2672677467393709,\quad 4108255612276161,\\
2806739366867294,\quad 719107361052288\\
\end{gathered}
\]

\paragraph{Reference cubic margin.} After $\nu=z+1$, the numerator proving $N_m^{\mathrm{ref}}>1/100$ has coefficients
\[
\begin{gathered}
3790502986637265684840,\quad 17978298402717728137711,\quad 33955060177057334483573,\\
31899843595051464379292,\quad 14895188986348368035728,\quad 2762610207555043720056\\
\end{gathered}
\]

\paragraph{Gauge upper-bound tail.} For $\nu\ge2$, the polynomial upper bound used in $\tau_m(L)<1/20$ has coefficients
\[
\begin{gathered}
-189709065/2,\quad -507201030,\quad -935658,\\
58481\\
\end{gathered}
\]

\paragraph{Remaining scalar comparisons.}
The certificate also checks exactly
\[
 \frac{1760850}{91}-10253>0,
 \qquad
 \frac4{462105}\left(\frac{1760850}{90}-10253\right)<\frac1{10},
\]
which give $-1/10<S_m<0$, and combines
$N_m^{\rm ref}>1/100$, $S_m>-1/10$, and
$\tau_m(L)<1/20$ to obtain $N_m(L)>1/200$.


\section{Near-threshold control example}

For the affine path
\[
p=\varepsilon,\quad
u=1+(2-t)\varepsilon+\theta\varepsilon^2,\quad
v=t\varepsilon-\theta\varepsilon^2,
\]
\[
q=\frac12-\left(\frac12+t\right)\varepsilon
 +\left(\theta-\frac M2\right)\varepsilon^2,
\]
the critical vector tends to one half of the homogeneous kernel vector.  The
diffusion excess over the exact linear threshold is
\[
 d_Z-8\sum_{j=2}^{m-1}d_j
 =4\left(M+6t+3t^2-\frac{t^2}{m-2}\right)\varepsilon^2
 +O(\varepsilon^3).
\]
For $(t,\theta,M)=(2/9,1/2,1)$ and $m=3$, exact elimination yields
\[
 c_3(\varepsilon)=\frac6{1379}
 +\frac{421985}{11409846}\varepsilon+O(\varepsilon^2),
\]
and rational remainder bounds prove $c_3(\varepsilon)>0$ on
$0<\varepsilon\le10^{-3}$.  This is a rigorous subcritical control example,
not a universal nonlinear lower bound.

\section{Semilinear stability and robustness}

On the real fixed-integrated-mass $L^2$ space, the positive diagonal Neumann
diffusion operator has domain $H_N^2$, is sectorial, and generates an analytic
semigroup.  Its half-order fractional domain is $H^1$.  In one space
dimension, $H^1$ is a Banach algebra and embeds into $C^0$, so the quadratic
mass-action Nemytskii map is smooth from $H^1$ to $L^2$.  The patterned branch
lies in $H_N^2$.  The fixed-mass restriction removes the homogeneous
conservation zero, and Neumann boundary conditions on an interval leave no
continuous translation-neutral mode.  Strict spectral negativity therefore
implies local exponential asymptotic stability in $H^1$.

For robustness, perturb positive steady fluxes and equilibrium coordinates
within the positive-equilibrium realization manifold, reconstructing the
rate constants accordingly; also allow small perturbations of diffusion
ratios and interval length.  Use one scalar diffusion multiplier as the
implicit variable.  The critical eigenvalue derivative is nonzero, so the
implicit function theorem continues the crossing.  Low-mode gaps persist by
analytic perturbation.  A positive diffusion lower bound and a
reaction-Jacobian upper bound give one neighborhood controlling all high
modes.  Smooth dependence of the reduced normal-form coefficients preserves the
negative cubic sign and the branch spectral gap.  This is a retuned
codimension-one statement, not a claim that arbitrary nearby parameter points
are already critical, and no dimension-uniform radius is claimed.

\section{Numerical protocol and independent verification}

Cosine--Galerkin simulations use exact modal convolution for the quadratic
reaction terms and an adaptive stiff integrator.  Initial data are small
perturbations of $\1$ in the critical direction.  The first-mode amplitude is
measured by the adjoint projection.  Spatial and temporal refinement tables,
open-format profile data, exact parameters, and all solver tolerances are
included in the data archive.  No numerical output is used in a proof.

The release contains independent programs for the reaction family, flux cone,
all-spectrum block theorem, omission minors, principal-minor diffusion-ray theorem, diffusion
criterion, contrast bounds, improved profile, Pareto family, mode
certificates, harmonic corrections, cubic sign, and branch stability.
Mutation tests change signs, endpoints, the factor eight, Fourier factors,
and certificate coefficients and must be rejected.

\end{document}
